\chapter{Formázási tudnivalók}

Ez a fejezet összeszedi azokat a nélkülözhetetlen elemeket, amelyeket a dolgozat
\LaTeX-ben való elkészítése során használnod kell. Feltétlenül olvasd egyszer
végig, hátha akad köztük újdonság.

\section{Általános tudnivalók}

A diplomaterv szabványos méretű A4-es lapokra kerüljön. Az alapértelmezett
betűkészlet a 12 pontos Times New Roman, 1.15-ös sorközzel. Ez a sablon
(\texttt{sapientia-thesis.cls}) a fenti előírásoknak megfelelően van
kialakítva, így általában külön nem kell foglalkoznod velük.

Minden oldalon -- az első négy szerkezeti elem kivételével -- szerepelnie kell
az oldalszámnak.

A fejezeteket decimális beosztással kell ellátni. Az ábrákat a megfelelő helyre
be kell illeszteni, fejezetenként decimális számmal és kifejező címmel kell
ellátni -- ezt a \LaTeX\ automatikusan kezeli a \verb|\caption| és
\verb|\label| parancsok segítségével.

Az irodalomjegyzék szövegközi hivatkozása sorszámozva történik (\texttt{[1]},
\texttt{[2]} stb.). A teljes lista a hivatkozási sorrendben a szöveg végén
szerepel. A szakirodalmi források címeit mindig az eredeti nyelven kell
megadni. A listában szereplő valamennyi publikációra hivatkozni kell a
szövegben.

\section{Stílusok}

A \LaTeX\ szövegek egységességét környezetek (\textit{environments}) és
parancsok (\textit{commands}) segítségével lehet a legegyszerűbben garantálni.
A sablon előre definiált parancsokat tartalmaz a leggyakoribb esetekre
(fejezetek, ábrák, kódblokkok, hivatkozások). Azt javaslom, hogy a sablon
struktúráját kövesd, és ne definiálj saját formázó parancsokat -- ha valami
hiányzik, szólj a konzulensednek.

\section{Címsorok}

A fejezetcímek esetén a \verb|\chapter|, \verb|\section|, \verb|\subsection|
és \verb|\subsubsection| parancsokat használjuk. Negyedik szintnél mélyebbre
egy ilyen terjedelmű munkában ritkán van szükség; ha ez mégis felmerülne,
sokszor inkább a fejezetszerkezetet érdemes átgondolni újból.

Kerüljük azt, amikor egy fejezetcím után rögtön alfejezet következik -- legalább
pár sor bevezető szöveg legyen közöttük. Hasonlóan érdemes elkerülni azt is,
amikor a címsor után rögtön felsorolás következik.

\section{Másolás, beillesztés}

A copy-paste a szép formázás legnagyobb gyilkosa. Amennyiben külső forrásból
(weboldal, PDF, más szerkesztő) illeszted be a szöveget egy \texttt{.tex}
fájlba, figyelj arra, hogy az idézőjelek (\verb|„|~és~\verb|"|), kötőjelek
(\verb|--|, \verb|---|) és az ékezetes karakterek helyesen kerüljenek át.
\LaTeX-ben a speciális karaktereket (\verb|%|, \verb|&|, \verb|$|, \verb|#|
stb.) visszaperrel kell escape-elni: \verb|\%|, \verb|\&|, \verb|\$| stb.

Szöveg áthelyezésekor a kereszthivatkozások (\verb|\ref|, \verb|\cite|)
automatikusan érvényesek maradnak -- ez a \LaTeX\ egyik nagy előnye a
szövegszerkesztőkkel szemben.

\section{Fordítás és frissítés}
\label{sec:forditas}

A \LaTeX-ben a tartalomjegyzék, ábrajegyzék, kereszthivatkozások és
irodalomjegyzék automatikusan frissülnek fordításkor -- nincs szükség kézi
mezőfrissítésre. Hivatkozások változásakor azonban a teljes fordítási
sorozatot le kell futtatni:

\begin{enumerate}[noitemsep]
  \item \texttt{pdflatex main} -- első átmenet, aux-fájlok generálása
  \item \texttt{biber main} -- irodalomjegyzék feldolgozása
  \item \texttt{pdflatex main} -- hivatkozások beírása
  \item \texttt{pdflatex main} -- tartalomjegyzék véglegesítése
\end{enumerate}

A legtöbb \LaTeX-szerkesztő (TeXstudio, VS~Code + LaTeX Workshop, Overleaf)
ezt a sorozatot egyetlen gombnyomásra elvégzi. Ha csak a szöveg változott és
nincs új hivatkozás, elegendő egyszer fordítani.

\section{Helyesírás}

A rossz helyesírásra nincs mentség. E fejezetben összeszedem a leggyakrabban
látott hibákat, amiknek elkerülésére érdemes odafigyelni. Ettől függetlenül
melegen ajánlom, hogy a kész dolgozatod olvastasd át egy
barátoddal/családtagoddal, hogy az apróbb, megbúvó hibákat is kiszűrd.

\subsection{Elgépelések}

Ez mindenkivel megesik. A TeXstudio és a VS~Code~+~LaTeX~Workshop bővítmény
beépített helyesírás-ellenőrzővel rendelkezik; győzödj meg arról, hogy a
\textbf{magyar szótár} telepítve és beállítva van. Overleaf-en a
\textit{Spell check} menüben válaszd a \textit{Hungarian} nyelvet.

\subsection{Egyeztetés hiánya}

Az elírások egyik leggyakoribb formája az egyes szám/többes szám
egyeztetésének hiánya mondatrészek között, mint például: „Petike és a
barátnője elmentek a boltba és hozott egy kiló kenyeret." Ezek a mondatok
főleg az utólagos átfogalmazások, belejavítások során keletkeznek; legjobb
védelem ellenük az utólagos átolvasás.

\subsection{Külföldi szavak, kifejezések}

Az idegen szavakkal csak a baj van. Az általános jó tanácsom, hogy amennyiben
csak lehetséges, \textbf{magyar, vagy magyarosított írásmódú szakkifejezéseket
használj}, és a könnyebb olvashatóság érdekében mindig kerüld az idegen szavak
ragozását. Pl. „Apache-csal" helyett írjuk azt, hogy „Apache webszerverrel".

Néhány gyorstipp szoftverfejlesztőknek: property $\to$ tulajdonság; event $\to$
esemény; method $\to$ metódus/függvény; debug $\to$ hibakeresés; file $\to$
fájl.

\subsection{Helyesírás-ellenőrző}

Személyes ízlés kérdése, hogy milyen szerkesztővel dolgozol, ugyanakkor azt
\textbf{meg kell oldanod, hogy legyen meg mellé magyar nyelvű
helyesírás-ellenőrződ}. E nélkül dokumentumot szerkeszteni olyan, mint
papíron programozni.

TeXstudio esetén: \textit{Options / Configure TeXstudio / Editor / Spelling
dictionary} alatt telepítse a \texttt{hu\_HU} szótárat. Tesztelésképpen írd
be, hogy \textit{„elgemépelés"} -- ha nem húzza alá pirossal, valami nincs
rendben a beállításaiddal.

\section{Képek}

A dolgozatodban valószínűleg számos ábrára lesz szükséged.

\subsection{Beszúrás, formázás}

A képhez használd a \texttt{figure} környezetet, a képaláírást a
\verb|\caption| paranccsal add meg, és adj neki egy \verb|\label|-t a
hivatkozáshoz. A sablonban az \texttt{[H]} elhelyezési opció (a \texttt{float}
csomag) rögzíti az ábrát a szöveghez képest, ami diplomadolgozatnál általában
a legjobb választás.

\begin{figure}[H]
  \centering
  \fbox{\rule{0pt}{4cm}\rule{0.65\textwidth}{0pt}}
  \caption{Példa képaláírásra}
  \label{fig:pelda}
\end{figure}

Az ábra sorszáma mellé mindig érdemes rövid magyarázatot is fűzni. Erre
láthatunk példát a~\ref{fig:pelda}. ábrán. Vektorgrafikus ábrákat PDF vagy
EPS formátumban illeszd be (\verb|\includegraphics{abra.pdf}|); ezek
nyomtatásban is élesen jelennek meg.

Tapasztalni fogod, hogy a \texttt{figure} és \texttt{table} lebegő elemek
néha nem a kívánt helyre kerülnek. Az ábrák méretének módosításával, vagy a
\verb|[H]| opció alkalmazásával tudod megoldani, hogy ne maradjanak nagy üres
felületek a dolgozatodban.

\subsection{Képminőség}

A raszteres (tehát nem vektorgrafikus) képek használata különös
körültekintést igényel. Ezek kiválóan néznek ki a 90 dpi-s monitorodon, ám a
600/1200~dpi-s nyomtatókon kinyomtatva rendkívül bénák lesznek a szép,
pixelmentes szövegek és vektorgrafikus ábrák mellett.

Ha tehát lehetséges, használjunk vektorgrafikus ábrákat (PDF, EPS, SVG), vagyis
a diagramokat, forráskódot stb. ne képernyőképeken keresztül, hanem közvetlen
\verb|\includegraphics| megoldással illesszük be. Ha elkerülhetetlen a
raszteres képek használata, próbáljunk meg minél magasabb felbontású képet
berakni.

\section{Kereszthivatkozások}

Amennyiben szeretnél egy ábrára, táblázatra, képletre vagy korábbi fejezetre
hivatkozni, helyezz el egy \verb|\label{azonosito}| parancsot a hivatkozott
elemnél, majd a szövegben hivatkozz rá a \verb|\ref{azonosito}| paranccsal.
Például: lásd a~\ref{fig:pelda}. ábrát, vagy lásd a~\ref{sec:forditas}.
fejezetet.

Kerüld az „előző oldalon", „következő oldalon" fordulatokat, ugyanis az ábrák
végső helyzete a tördelés során még megváltozhat.

\section{Irodalomhivatkozások}
\label{sec:irodalom}

Az irodalomhivatkozások kezelésére a \texttt{biblatex} csomagot és a
\texttt{biber} háttérprogramot használjuk. A forrásokat a
\texttt{references.bib} fájlban kell felvenni BibTeX formátumban. Amikor
dolgozatodban egy külső műre, weboldalra, könyvre, előadásra stb.\ szeretnél
hivatkozni, add hozzá a \texttt{.bib} fájlhoz, majd a szövegben hivatkozz rá a
\verb|\cite{kulcs}| paranccsal~\cite{Abb00, Del03}.

Folyóirat cikkeknél a szerzők mellett szerepeljen a pontos cím, a folyóirat
neve, évfolyam, szám és oldalszám. Internet hivatkozásoknál fontos, hogy az
oldal szerzőjét, címét és az elérés időpontját is add meg (mivel a link egy
idő után elérhetetlenné is válhat) -- a BibTeX \texttt{@online} típusával és
az \texttt{urldate} mezővel.

A sablon \texttt{numeric} stílusú hivatkozásokat használ: a szövegben a
forrásra \texttt{[1]}, \texttt{[2]} alakban hivatkozunk, a lista a szöveg
végén, hivatkozási sorrendben jelenik meg.

\subsection{Pozícionálás}

Az irodalomhivatkozások a szövegtörzsben, ábrák képaláírásaiban és
táblázatokban is előfordulhatnak, de \textbf{fejezetcímekben soha}. Ha egy
adott forrás egy egész bekezdésre vonatkozik, akkor is elég, ha az első mondat
vagy bekezdés után megemlítjük.

\subsection{Mikor kell hivatkoznom?}

Minden külső forrásból átvett képnél, szövegrésznél, illetve olyan
állítások, technológiák, algoritmusok megemlítésénél, melyek nem feltétlenül
egyértelműek egy átlagos műveltségű olvasó számára~\cite{Hop95}.
Nagyságrendileg egy szakdolgozatban átlagosan 10--20, egy diplomatervben
átlagosan 20--30 külső forrást illik megemlíteni. Ha lehet, törekedjél
nyomtatott forrásokra (folyóiratcikkek, könyvek, konferencia-kiadványok), és
csak akkor hivatkozz weboldalakra, ha nyomtatott forrást nem találtál.

\section{\LaTeX\ tippek és trükkök}

\subsection{Navigáció a dokumentumban}

Mivel a fejezeteket nagy valószínűséggel nem sorrendben fogod tartalommal
feltölteni, érdemes kihasználni a szerkesztő navigációs lehetőségeit. A
TeXstudio bal oldali paneljén a \textit{Structure} nézet mutatja az összes
fejezetet és alfejezetet, egy kattintással az adott helyre ugrik. A
\texttt{hyperref} csomag (a sablonban már betöltve) a generált PDF-be kattintható
könyvjelzőket és belső linkeket szúr be, amelyek PDF-megjelenítőkben
navigációs panelként jelennek meg.

\subsection{Megjegyzések a forráskódban}

A \LaTeX-ban a megjegyzéseket a \verb|%| karakterrel lehet beilleszteni -- a
sor többi részét a fordító figyelmen kívül hagyja. Ezt előszeretettel
használhatod, hogy megjelöld azokat a részeket, ahova még vissza kell térned.
Konzulensed a \texttt{changes} csomag segítségével tud nyomtatható
megjegyzéseket és nyomonkövethető javításokat fűzni a dokumentumhoz.

\subsection{Korrektúra és együttműködés}

Konzulensed legegyszerűbben a \texttt{changes} csomag
(\verb|\added{...}|, \verb|\deleted{...}|, \verb|\replaced{...}{...}|)
segítségével tud megjegyzéseket és javításokat fűzni a dokumentumhoz, amelyek
a PDF-ben jól látható módon jelennek meg. Overleaf esetén a beépített Track
Changes funkciót is használhatjátok.

\subsection{Hasznos gyorsbillentyűk -- TeXstudio}

A TeXstudio számos beépített gyorsbillentyűt kínál. A legfontosabbak:
\texttt{F5} -- fordítás és megjelenítés; \texttt{F6} -- csak fordítás;
\texttt{Ctrl+Shift+B} -- \texttt{biber} futtatása; \texttt{Ctrl+Klikk} -- ugrás
a forrásból a PDF-be (és vissza). Tetszőleges parancshoz saját
gyorsbillentyű rendelhető az \textit{Options / Configure TeXstudio / Shortcuts}
menüben.

\section{Kódrészletek}

Érdekesebb és bonyolultabb programozási megoldásainkat bátran illusztrálhatjuk
kódrészletek beszúrásával. Fontos, hogy a beillesztett kódrészlet mérete álljon
arányban annak fontosságával -- ritkán érdemes egy bekezdésnyi kódnál többet
beszúrni egyszerre. Amennyiben egy bonyolultabb, akár több oldalas algoritmust
szeretnénk bemutatni, annak kódját érdemesebb függelékbe rakni.

\subsection{Formázás}

A kódrészletek beillesztésére a \texttt{listings} csomagot használjuk
(\texttt{lstlisting} környezet). A sablonban a csomag alapértelmezetten
kulcsszókiemelést és monospace betűtípust alkalmaz. Kódrészletek
beillesztésére semmiképpen ne képeket használjunk.

A kódrészletek formázásánál kerüljük a helypazarlást, illetve próbáljunk
megelőzni az olvashatóságot rontó sortördelést, akár a forráskód módosításának
árán is. Másolás előtt érdemes a behúzások mértékét 4-ről 2 karakterre
csökkenteni (Visual Studio-ban: \textit{Tools / Options / Text Editor / C\# /
Tabs}: Tab size = 2, Indent size = 2).

Az alábbi C\# kódrészlet szemlélteti a \texttt{lstlisting} környezet
használatát:

\begin{lstlisting}[language={[Sharp]C}]
static void Main(string[] args)
{
  var ci = new CultureInfo("en-us");
  ci.NumberFormat.CurrencySymbol = "";
  Thread.CurrentThread.CurrentCulture = ci;
  Console.WriteLine(ci);
  Console.WriteLine("{0:c}", 5.66);
}
\end{lstlisting}

A programozási elemek (osztálynevek, függvények) szövegközi kiemelésére
használd a \verb|\texttt{NévMinták}| parancsot; ez egységes, monospace
megjelenést biztosít.

\subsection{Irodalomjegyzék}

Az irodalomjegyzékben szereplő hivatkozásokat a \texttt{biblatex} csomag kezeli.
A \verb|\printbibliography| parancs automatikusan generálja a listát a
\texttt{references.bib} fájl tartalma alapján. A fordítási sorozatban a
\texttt{biber} futtatása frissíti az irodalomjegyzéket (lásd a~\ref{sec:forditas}.
fejezetet). Csak azok a tételek jelennek meg a listában, amelyekre a szövegben
\verb|\cite|-tal hivatkoztál -- a \verb|\nocite{*}| paranccsal az összes tétel
megjeleníthető.

\section{Utolsó simítások}

Miután elkészültünk a dokumentációval, ne felejtsük el a következő lépéseket:

\begin{itemize}[noitemsep]
  \item \textbf{Teljes fordítási sorozat:} \texttt{pdflatex} $\to$
        \texttt{biber} $\to$ \texttt{pdflatex} $\to$ \texttt{pdflatex}.
        Ellenőrizd, hogy nem jelent-e meg valahol az „undefined reference"
        vagy „Citation ... undefined" figyelmeztetés a logban.
  \item \textbf{PDF metaadatok:} A \verb|\hypersetup| paranccsal a
        \texttt{pdfauthor}, \texttt{pdftitle} és \texttt{pdfkeywords} mezők
        kitölthetők -- ezek a PDF-fájl tulajdonságaiban jelennek meg.
  \item \textbf{Kinézet ellenőrzése PDF-ben:} A legjobb teszt a végén, ha a
        kész PDF-et átolvasod, különös tekintettel az oldalelrendezésre, az
        ábrák pozíciójára, a tartalomjegyzék oldalszámaira és az
        irodalomjegyzékre.
\end{itemize}
